% ============================================================================
% sections/abstract.tex
% ============================================================================
\begin{center}
{\Large\bfseries\color{navy} Abstract}
\end{center}
\addcontentsline{toc}{section}{Abstract}
\vspace{0.4em}

Retrieval-Augmented Generation (RAG) systems answer questions by retrieving
supporting documents from a corpus and grounding a generated answer in
them. This architecture introduces a critical attack surface: an adversary
who can influence either the corpus or the query can steer the system
toward a false answer without touching the model's weights. This project
surveys six frontier papers (2025--2026) documenting increasingly
effective poisoning attacks --- compound query-and-corpus attacks
(PIDP-Attack), semantically fluent detection-evading poison
(SilentRetrieval), and coordinated multi-document collusion --- alongside
a sequence of defenses: spectral pre-ingestion filters (DRS), query-time
consensus reshaping (ShieldRAG / ``Push and Pull''), adversarially trained
retrievers with counterfactual verification (RAGuard / ZKIP), and
multi-ring defense-in-depth pipelines (TriShieldRAG). Read together, these
works reveal that no single existing defense is sufficient: spectral
filters lack query-time protection, single-document counterfactual
verification is blind to colluding poison documents, and multi-ring
pipelines can still collapse into false consensus when every stage
evaluates the same already-poisoned retrieved context.

\vspace{0.6em}
This project implements and empirically evaluates OmniGuard-RAG, a
four-ring defense-in-depth framework built to close these gaps: (1) a
query-path guard that strips adversarial suffixes before retrieval, (2) a
spectral ingestion filter with a fit/calibration-split PCA to avoid
self-referential thresholds, (3) a risk-aware router using two independent
signals (embedding cohesion and answer-vote contention), and (4) a
Group-Wise Counterfactual Consensus (GWCC) mechanism that generalises
leave-one-out verification to leave-group-out exclusion for colluding
document cliques. A persistent, cross-query Dynamic Trust Store
accumulates evidence about individual documents across a session. Every
attack and every defense operates on real, generated text embedded through
a real, fixed TF-IDF vector space --- no detection threshold is hand-set;
every one is calibrated from measured clean-data statistics.

\vspace{0.6em}
Evaluated across 8 independently regenerated seeds (1,600 queries in
total) against five baseline defenses under six attack regimes (clean,
standard poisoning, PIDP, collusion, stealth collusion, and
silent/SilentRetrieval-style poisoning), OmniGuard-RAG reaches
$100.0\pm0.0\%$ accuracy and $0.0\pm0.0\%$ overall attack success rate,
with stealth-collusion attack success at a small and honestly-reported
$0.1\pm0.1\%$ --- not a fabricated perfect zero. A per-ring ablation shows
that no single ring achieves this result alone: Group-Wise Counterfactual
Consensus has a real, measured ceiling of roughly 8--10\% attack success
against camouflaged collusion within a single query, and the headline
result is reached only because the persistent trust store closes the gap
that a single query cannot structurally close. Eight real, measured
software bugs --- including a self-referential filter threshold and a
consensus mechanism that was silently a no-op --- were found and fixed
during development through direct measurement rather than parameter
tuning, and this debugging trail is documented as part of the project's
methodology, not hidden from it.

\vspace{0.8em}
\noindent\textbf{\color{navy}Keywords:} \textit{Retrieval-Augmented
Generation, Data Poisoning, Defense-in-Depth, Counterfactual Consensus,
Spectral Filtering, Adversarial Robustness, TF-IDF Retrieval.}
